% ============================================================
%  Pre-registro del protocolo experimental - TFG AutoPipe
%  Autor: Ricardo Antonio Borbon Mena
%
%  PROPOSITO
%  ---------
%  Documento fechado y versionado que fija, ANTES de ejecutar los
%  experimentos, las decisiones de protocolo que el anteproyecto
%  (seccion 4.3) no especifica. No modifica ningun criterio de
%  exito ya comprometido: solo cierra sus silencios.
%
%  REGLA DE ORO
%  ------------
%  Ninguna version de este documento se escribe DESPUES de la
%  ejecucion que gobierna:
%    v0.1  semana 1  - decisiones de protocolo (este documento)
%    v0.2  semana 5  - antes de correr el cribado de sensibilidad (cerrada)
%    v1.0  semana 9  - antes de ejecutar el leave-one-dataset-out
%
%  Cada version se cierra con un commit propio y una etiqueta
%  anotada de git (prereg-v0.1, prereg-v0.2, prereg-v1.0).
%
%  Compilar con:  pdflatex Preregistro.tex
% ============================================================

\documentclass[11pt,letterpaper]{article}

\usepackage[utf8]{inputenc}
\usepackage[T1]{fontenc}
\usepackage[spanish,es-tabla]{babel}
\usepackage{lmodern}
\usepackage[letterpaper,top=2.4cm,bottom=2.4cm,left=2.8cm,right=2.5cm]{geometry}
\usepackage{xcolor}
\usepackage{colortbl}
\usepackage{booktabs}
\usepackage{array}
\usepackage{tabularx}
\usepackage{enumitem}
\usepackage{fancyhdr}
\usepackage{titlesec}

\definecolor{tecnavy}{HTML}{003366}
\definecolor{gris}{HTML}{666666}
\definecolor{tablehead}{HTML}{E6EEF5}

\newcolumntype{Y}{>{\raggedright\arraybackslash}X}
\newcolumntype{L}[1]{>{\raggedright\arraybackslash}p{#1}}

\setlength{\parindent}{0pt}
\setlength{\parskip}{5pt}
\renewcommand{\arraystretch}{1.2}
\setlist[itemize]{leftmargin=1.4em,topsep=2pt,itemsep=3pt}
\setlist[enumerate]{leftmargin=1.6em,topsep=2pt,itemsep=3pt}

\titleformat{\section}{\large\bfseries\sffamily\color{tecnavy}}{\thesection}{0.6em}{}
\titleformat{\subsection}{\normalsize\bfseries\sffamily}{\thesubsection}{0.5em}{}

\pagestyle{fancy}
\fancyhf{}
\renewcommand{\headrulewidth}{0.4pt}
\fancyhead[L]{\small\textit{AutoPipe}}
\fancyhead[R]{\small Pre-registro \preregversion}
\fancyfoot[C]{\thepage}

% --- Identidad de la version --------------------------------
\newcommand{\preregversion}{v1.0}
\newcommand{\preregfecha}{7 de septiembre de 2026}
\newcommand{\preregtag}{prereg-v1.0}

% Fechas de cierre de las versiones anteriores, para el historial.
\newcommand{\preregfechauno}{21 de julio de 2026}
\newcommand{\preregfechados}{20 de agosto de 2026}

\begin{document}

% =========== PORTADA ========================================
\begin{center}
{\LARGE\bfseries\color{tecnavy} Pre-registro del protocolo experimental}\\[6pt]
{\large AutoPipe --- Trabajo Final de Graduación}\\[4pt]
{\small Ricardo Antonio Borbón Mena --- Instituto Tecnológico de Costa Rica}\\[14pt]

\begin{tabularx}{0.92\textwidth}{L{3.6cm} Y}
\toprule
\rowcolor{tablehead} \multicolumn{2}{l}{\textbf{Identificación de esta versión}} \\
\midrule
Versión & \preregversion \\
Fecha de cierre & \preregfecha \\
Estado & \textbf{Versión congelada previa a la ejecución del experimento.} Cierra
las dos decisiones que quedaban abiertas —el número de semillas y el tratamiento
de la región inevaluable— y registra las reglas del recomendador y de la variable
medida. Se cierra \textbf{antes} de ejecutar el
\textit{leave-one-dataset-out}: ningún resultado del experimento existía cuando
se escribió. \\
Etiqueta de git & \texttt{\preregtag}, etiqueta anotada que apunta al
\emph{commit} que cierra esta versión. El fechado verificable es la fecha de
\textbf{ese commit}, no la de creación de la etiqueta, y se obtiene con
\texttt{git log -1 \preregtag}. Véase la fe de erratas de la
Sección~\ref{sec:erratas}. \\
Forma de cita en el informe & «Pre-registro \preregversion, etiqueta
\texttt{\preregtag}, \preregfecha» \\
\bottomrule
\end{tabularx}

\vspace{6pt}
{\small\color{gris} El documento no lleva escrito dentro el hash del commit que lo
congela: ese hash se calcula sobre el contenido del documento, de modo que
escribirlo dentro de él lo alteraría. Por la misma razón tampoco lleva la fecha de
ese commit. Ambos se obtienen desde fuera, con \texttt{git log -1 \preregtag}.\par}
\end{center}

\vspace{10pt}

% =========== 1. PROPOSITO ===================================
\section{Propósito y alcance de este documento}

Este documento fija, con fecha anterior a la ejecución de los experimentos, las
decisiones de protocolo que el anteproyecto no especifica. Su función es que las
reglas de análisis no puedan ajustarse después de conocer los resultados.

\textbf{Este documento no modifica ningún criterio de éxito del anteproyecto.} La
Sección~4.3 del anteproyecto (corpus y presupuesto, métricas primarias, protocolo
de comparación, criterio de éxito de la hipótesis y criterios de éxito de
ingeniería) permanece vigente sin cambios. Lo que aquí se registra son
precisiones sobre puntos en los que ese texto guarda silencio, más decisiones de
diseño internas que el texto no gobierna. La Tabla~\ref{tab:relacion} deslinda
ambas cosas explícitamente.

\begin{table}[htbp]
\centering
\caption{Deslinde entre lo que fija el anteproyecto y lo que precisa este documento.}
\label{tab:relacion}
\small
\begin{tabularx}{\textwidth}{L{4.1cm} L{4.3cm} Y}
\toprule
\rowcolor{tablehead} \textbf{Asunto} & \textbf{Lo que fija el anteproyecto \S4.3} & \textbf{Lo que precisa este documento} \\
\midrule
Modelo base &
«Un clasificador sensible al preprocesamiento (regresión logística o $k$ vecinos
más cercanos de scikit-learn)», fijo en ambas condiciones. &
Se elige $k$ vecinos más cercanos y se congela su configuración completa
(Sección~\ref{sec:modelo}). El texto autoriza expresamente esta opción. \\
\midrule
Espacio de búsqueda &
Nada: \S4.3 no lo especifica. &
Parametrización relativa (fraccional) y espacio canónico versionado, con su
composición congelada (Sección~\ref{sec:espacio}). \\
\midrule
Meta-características &
Nada: \S4.3 no lo especifica. &
Biblioteca y grupos, submuestreo del grupo de complejidad, imputación previa y
regla de descarte de valores no finitos (Sección~\ref{sec:metacar}). \\
\midrule
Corpus &
Criterios de inclusión: licencia abierta, 500--50\,000 instancias, $\leq 100$
características, sin imagen ni texto embebidos, $\leq 10$ clases. Entre 15 y 25
datasets. 50 ensayos por dataset. &
Cómo se eligen los 15--25 entre todos los elegibles: desempate por sensibilidad
al preprocesamiento (Sección~\ref{sec:corpus}). Los criterios de inclusión no se
alteran ni se añaden nuevos. \\
\midrule
Semillas &
«El mismo presupuesto de 30 ensayos y las mismas semillas» en ambas condiciones. &
Cuántas semillas, cuáles literalmente, y cómo se agregan
(Sección~\ref{sec:lodo}). \\
\midrule
Ensayos hasta el objetivo &
«Número de ensayos necesarios para alcanzar al menos el 95\% del mejor valor de
métrica hallado por el arranque en frío con el presupuesto completo». &
Respecto a qué corrida en frío se define el objetivo, si los ensayos inyectados
cuentan, y qué se hace cuando una condición no alcanza el objetivo
(Sección~\ref{sec:lodo}). \\
\midrule
Criterio de éxito &
Reducción $\geq 30\%$ de la mediana de ensayos hasta el objetivo, con
significancia según Wilcoxon pareada por dataset, $\alpha = 0{,}05$. &
Lateralidad de la prueba y análisis secundarios, que no sustituyen al primario
(Sección~\ref{sec:analisis}). \\
\bottomrule
\end{tabularx}
\end{table}

% =========== 2. MODELO BASE =================================
\section{Modelo base congelado}
\label{sec:modelo}

El modelo base es \texttt{KNeighborsClassifier} de scikit-learn, con esta
configuración literal, idéntica en todas las condiciones, datasets y semillas:

% \shorthandoff{"} es imprescindible aquí: con babel en español, la comilla es un
% carácter activo y la secuencia "u compone «ü», de modo que weights="uniform" se
% imprimía como weights=üniform". Afecta justo al literal que hay que copiar.
\begin{center}
\shorthandoff{"}%
\texttt{KNeighborsClassifier(n\_neighbors=5, weights="uniform",\\
metric="minkowski", p=2, algorithm="brute")}%
\shorthandon{"}
\end{center}

\textbf{Justificación de la elección.} La Sección~4.3.3 del anteproyecto autoriza
expresamente regresión logística o $k$ vecinos más cercanos. Se elige el segundo
porque opera sobre distancias, de modo que el escalado, la imputación y la
selección de características afectan la métrica de forma sustantiva. Un modelo
base insensible al preprocesamiento produciría un efecto piso: si la métrica
apenas varía entre pipelines, el objetivo del 95\% se alcanza en el primer
ensayo de ambas condiciones y la comparación queda vacía por construcción,
independientemente de la calidad del recomendador. La elección de $k$ vecinos más
cercanos es la mitigación más económica de ese riesgo.

\textbf{Justificación de} \texttt{algorithm="brute"}\textbf{.} Los algoritmos
basados en árboles degradan su ventaja a un centenar de dimensiones y su
construcción puede introducir variación entre ejecuciones. Forzar fuerza bruta
elimina esa fuente de variabilidad y sirve al criterio de reproducibilidad de la
Sección~4.3.5 del anteproyecto (la misma semilla debe producir el mismo pipeline
ganador).

\textbf{Consecuencia sobre el costo.} El modelo invierte el perfil de cómputo:
entrenar es trivial y predecir crece cuadráticamente con el número de instancias,
de modo que el costo total queda dominado por los datasets cercanos al techo de
50\,000 instancias. La medición de ese costo se realiza en la semana 6 y se
registra en la versión v1.0 de este documento, que es la primera posterior a esa
medición.

Ningún hiperparámetro del modelo base entra en el espacio de búsqueda: el
anteproyecto exige un modelo base fijo para que las diferencias observadas sean
atribuibles al pipeline de preprocesamiento.

% =========== 3. ESPACIO DE BUSQUEDA =========================
\section{Espacio de búsqueda}
\label{sec:espacio}

\subsection{Parametrización relativa}

Todo parámetro cuyo dominio dependa del dataset se expresa en términos
\textbf{relativos} y no absolutos. En particular, la selección de
características se parametriza como fracción de columnas conservadas, no como un
número absoluto de columnas.

\textbf{Razón.} El arranque en caliente transfiere configuraciones entre
datasets. Un parámetro absoluto no es transferible: conservar diez columnas
significa algo distinto en un dataset de doce características que en uno de cien,
y el valor podría ser inválido en el dataset receptor. La parametrización
relativa es condición necesaria para que una configuración recomendada sea
ejecutable en un dataset que el recomendador nunca vio.

\subsection{Espacio canónico y versionado}

El espacio de búsqueda tabular es único, idéntico para todos los datasets del
corpus, y se identifica con un número de versión que se registra junto a cada
entrada de la base de meta-conocimiento y junto a cada resultado experimental.

\textbf{Razón.} En la implementación del motor de búsqueda, los parámetros que
Optuna registra por ensayo son \emph{índices enteros} sobre la lista de opciones
de cada paso, más un índice adicional que representa omitir el paso. Un índice no
tiene significado por sí mismo: depende del orden de la lista. Si el espacio
cambiara de claves o de orden entre datasets, una configuración inyectada como
arranque en caliente designaría transformaciones distintas de las previstas, sin
producir ningún error visible. El versionado hace detectable esa condición:
una recomendación proveniente de una versión del espacio distinta de la vigente
se rechaza en lugar de ejecutarse.

\subsection{Composición congelada}

El espacio vigente se identifica como \texttt{tabular-v1}. Sus cuatro pasos, en
este orden de aplicación, son:

\begin{center}
\texttt{imputation} $\rightarrow$ \texttt{scaling} $\rightarrow$
\texttt{encoding} $\rightarrow$ \texttt{selection}
\end{center}

\begin{center}
\begin{tabularx}{0.92\textwidth}{L{2.6cm} Y r}
\toprule
\rowcolor{tablehead} \textbf{Paso} & \textbf{Opciones, en orden de índice} & \textbf{Con omisión} \\
\midrule
\texttt{imputation} & media, mediana & 3 \\
\texttt{scaling} & estándar, mín-máx, robusto & 4 \\
\texttt{encoding} & \textit{one-hot}, ordinal & 3 \\
\texttt{selection} & fracción 0,25 / 0,50 / 0,75, criterio \texttt{f\_classif} & 4 \\
\bottomrule
\end{tabularx}
\end{center}

Cada paso recibe además, de forma automática, la opción de \textbf{omitirse}, de
modo que el espacio contiene $3 \times 4 \times 3 \times 4 = 144$ pipelines. La
tabla de índices exacta, generada desde el propio código, está en
\texttt{docs/experimentos/ESPACIO-tabular-v1.md}.

\textbf{El escalado precede a la codificación} para que actúe solo sobre las
columnas numéricas originales y deje las indicadoras \textit{one-hot} en 0/1.
Estandarizarlas llevaría la indicadora de una categoría rara a valores cercanos a
$\pm 10$ y, con el modelo base de la Sección~\ref{sec:modelo}, esa sola columna
dominaría la distancia euclídea: el resultado mediría la rareza de una categoría y
no la calidad del preprocesamiento.

\textbf{Cambiar cualquier elemento de esa composición obliga a incrementar el
identificador de versión}, porque las configuraciones ya guardadas dejarían de ser
interpretables.

\subsection{Región inevaluable del espacio}

Los motores generan candidatos con pasos omitidos, y algunos son inevaluables:
omitir la codificación en un dataset con columnas categóricas deja una tabla que
el modelo no puede consumir, y omitir la imputación deja valores nulos. Esos
candidatos se registran con métrica \texttt{NaN} y nunca resultan ganadores.

La fracción medida sobre datos sintéticos va de 0 (dataset numérico sin
faltantes) a 56\% (con categóricas y con faltantes). Tiene dos consecuencias sobre
el protocolo:

\begin{enumerate}
  \item Sobre el cribado de sensibilidad, resuelta en la
        Sección~\ref{sec:corpus}.
  \item Sobre la variable «ensayos hasta el objetivo»: el arranque en caliente
        inyecta configuraciones ganadoras, que por construcción no caen en la
        región inevaluable, mientras que el frío gasta ahí parte de sus ensayos
        aleatorios. Se resuelve a continuación.
\end{enumerate}

\textbf{Decisión cerrada en esta versión.} La región inevaluable \textbf{se deja
como está}: ni el espacio ni los motores se modifican, y los candidatos que no se
pueden evaluar siguen registrándose con métrica \texttt{NaN} y consumiendo su
ensayo. En su lugar se declara una covariable, descrita en la
Sección~\ref{sec:covariables}.

La alternativa que la v0.2 contemplaba era introducir la noción de \emph{paso
obligatorio} en los motores, de modo que no se pudiera omitir la codificación en
un dataset con categóricas. Se descarta por tres razones, y ninguna es de
comodidad. Primera: sería un cambio de comportamiento público de los motores
introducido después de que el espacio quedara congelado, lo que obligaría a
incrementar el identificador de versión. Segunda: eso invalidaría la base de
meta-conocimiento ya construida, cuyos índices solo significan algo respecto de
\texttt{tabular-v1}, y habría que reconstruirla entera. Tercera, y la que más
pesa: cambiaría el espacio sobre el que se mide justo antes de medir, y el efecto
de ese cambio sobre la variable de interés sería favorable al arranque en
caliente, porque eliminaría precisamente los ensayos que el frío desperdicia. Una
modificación que agranda el efecto esperado, decidida después de conocer la
magnitud del problema, es exactamente lo que un pre-registro existe para impedir.

Se asume por tanto que la asimetría existe, se declara aquí, y se controla por
dos vías independientes: la covariable de la Sección~\ref{sec:covariables}, que la
hace visible dataset por dataset, y la ablación aleatoria de la
Sección~\ref{sec:analisis}, cuyas configuraciones provienen de datasets distintos
y por tanto sí pueden caer en la región inevaluable del dataset receptor.

% =========== 4. META-CARACTERISTICAS ========================
\section{Extracción de meta-características}
\label{sec:metacar}

El vector de meta-características es la entrada del recomendador: de él depende
qué datasets se consideran parecidos entre sí. Las reglas siguientes se fijan
aquí, antes de construir la base de meta-conocimiento, y salen de las mediciones
del spike B (\texttt{docs/experimentos/SPIKE-B-pymfe.md}).

\subsection{Biblioteca y grupos}

Se emplea \texttt{pymfe} 0.4.4 con cuatro grupos: \texttt{general},
\texttt{statistical}, \texttt{info-theory} y \texttt{complexity}. Producen 107
meta-características, y el conjunto de nombres es el mismo para cualquier
dataset, que es la condición para que el recomendador compare vectores alineados.

\subsection{Submuestreo del grupo de complejidad}

Los tres primeros grupos se extraen del \textbf{dataset completo}. El grupo
\texttt{complexity} se extrae de una \textbf{submuestra aleatoria de 500
instancias, estratificada por clase, con semilla 42}.

\textbf{Razón del submuestreo.} El costo de \texttt{complexity} crece de forma
cuadrática: el exponente medido es 1,91, y en el techo de 50\,000 instancias que
admite el anteproyecto la extracción costaría unos 21,6 minutos por dataset,
frente a menos de tres segundos para los otros tres grupos juntos.

\textbf{Razón del tamaño fijo, y no de una fracción fija.} Las medidas de
complejidad son propiedades geométricas de la muestra concreta, no estimadores de
una cantidad poblacional, de modo que cambian con el tamaño de la muestra. Dos
datasets solo son comparables si se les extrajo con el mismo tamaño. Como los
criterios de inclusión admiten datasets desde 500 instancias, \textbf{500 es el
único tamaño que todo dataset elegible puede alcanzar}. Un tamaño variable
cambiaría ruido por sesgo, y un sesgo correlacionado con el número de filas haría
que el recomendador juzgue similares a datasets que solo comparten su tamaño.

Si la estratificación no fuera posible por tener alguna clase menos miembros que
los que exige el reparto, se cae a un muestreo aleatorio simple con la misma
semilla y se deja constancia. Un dataset con menos de 500 instancias se usa
entero.

\subsection{Imputación previa}

Antes de invocar a \texttt{pymfe} se imputan los valores faltantes: mediana en las
columnas numéricas y moda en las categóricas. No es una comodidad: con valores
faltantes, el grupo \texttt{complexity} falla con una excepción de mensaje vacío
al precomputar el grafo de adyacencia.

\subsection{Determinismo}

Los valores se redondean a nueve decimales. Dos extracciones sobre el mismo
dataset difieren en \texttt{hubs.mean} y \texttt{hubs.sd} en torno a
$3 \times 10^{-15}$ por el orden de sumación en coma flotante; el redondeo
absorbe esa diferencia y vuelve el vector estable bit a bit, como exige el
criterio de reproducibilidad de la Sección~4.3.5 del anteproyecto.

\subsection{Meta-características no finitas}

Alrededor del 13\% del vector no es finito en un dataset dado, y \textbf{qué
posiciones lo son depende del dataset}: hay meta-características que existen para
unos y no para otros.

\textbf{Regla:} se descarta toda meta-característica que resulte no finita en
\textbf{algún} dataset del corpus. El vector final es, por tanto, la intersección
de las finitas. No se imputa ningún valor al nivel meta.

\textbf{Razón.} Imputar al nivel meta exigiría ajustar esa imputación por pliegue,
como la normalización, y añadiría una fuente de error difícil de auditar sobre una
métrica de similitud. Descartar es la operación conservadora.

El descarte se calcula \textbf{una sola vez}, sobre el corpus definitivo, y la
lista resultante se registra junto a la base de meta-conocimiento. Medido sobre
los 52 candidatos, 27 de las 107 meta-características son no finitas en al menos
uno, de modo que sobrevivirían 80; como el corpus definitivo es un subconjunto de
los candidatos, esa cifra es una \textbf{cota inferior} del número final.

% =========== 5. CORPUS ======================================
\section{Corpus y cribado de sensibilidad}
\label{sec:corpus}

Los criterios de inclusión son exclusivamente los de la Sección~4.3.1 del
anteproyecto. Este documento no añade, elimina ni relaja ninguno. Lo que fija es
el procedimiento para elegir los 15--25 datasets del corpus entre todos los que
resultan elegibles, que el anteproyecto no especifica.

\subsection{Procedimiento, fijado antes de ejecutarse}

\begin{enumerate}
  \item \textbf{Universo.} Todos los datasets que satisfacen los criterios de
        inclusión de la Sección~4.3.1, sin excepción.

  \item \textbf{Muestreo.} Doce pipelines por dataset candidato, extraídos
        uniformemente al azar del espacio de búsqueda canónico de la
        Sección~\ref{sec:espacio}, incluida la opción de omitir cada paso, con
        semilla fija \texttt{4242}.

        \textit{Pipelines inevaluables.} Parte del espacio no es evaluable sobre
        un dataset dado (Sección~\ref{sec:espacio}). Se \textbf{sigue extrayendo
        al azar, descartando los inevaluables, hasta reunir doce pipelines
        válidos}, y se registra cuántos se descartaron.

        \textit{Razón:} el cribado no solo mide, \emph{ordena} los candidatos por
        IQR y selecciona los primeros. Con doce extracciones fijas, el tamaño
        efectivo de muestra variaría entre doce y unos cinco según el dataset, de
        modo que el orden compararía estadísticas de precisión muy distinta.
        Remuestrear mantiene $n = 12$ en todos y hace comparables los IQR. Es
        además coherente con la razón por la que se eligió el IQR: mide si
        \emph{elegir entre} pipelines cambia el resultado, y un pipeline que no
        se puede ejecutar no es una de las opciones entre las que se elige.
        Contar los inevaluables como macro-F1 nula inflaría el IQR justo en los
        datasets con categóricas y faltantes, confundiendo «sensible al
        preprocesamiento» con «tiene columnas categóricas».

  \item \textbf{Evaluación.} El modelo base congelado de la
        Sección~\ref{sec:modelo}, la métrica primaria de clasificación del
        anteproyecto (macro-F1), y la misma partición única de entrenamiento y
        validación que emplea el protocolo de búsqueda.

  \item \textbf{Medida de sensibilidad.} El rango interquartílico (IQR) de los
        doce valores de macro-F1. Se registran además, con carácter descriptivo,
        el rango total ($\max - \min$) y el margen sobre el pipeline vacío
        (mejor pipeline aleatorio menos el desempeño sin preprocesamiento).

        \textit{Razón de la elección:} con doce muestras, el rango total lo
        determina un único pipeline degenerado. El IQR mide si \emph{elegir
        entre} pipelines cambia el resultado, que es precisamente lo que
        gobierna la variable «ensayos hasta el objetivo»; el margen sobre el
        pipeline vacío mide si el preprocesamiento \emph{ayuda}, que es una
        pregunta distinta y se reporta aparte.

  \item \textbf{Regla de familia.} Como máximo \textbf{un dataset por familia de
        origen}, aplicada \emph{después} del orden por IQR y \emph{antes} de las
        restricciones de diversidad. Distingue dos niveles:

        \textit{Nivel 1, mismas instancias (obligatorio).} Los cuatro
        \texttt{mfeat-*} son cuatro representaciones de las \textbf{mismas} 2\,000
        instancias. Si el corpus incluyera varias y la validación dejara una
        fuera, el recomendador seguiría teniendo el mismo dataset en su base bajo
        otra representación, y la independencia entre pliegues que supone el
        \textit{leave-one-dataset-out} dejaría de cumplirse. Es un problema de
        validez, no de diversidad.

        \textit{Nivel 2, misma fuente o esquema (recomendable).} Los seis del
        repositorio NASA MDP (\texttt{kc1}, \texttt{kc2}, \texttt{pc1},
        \texttt{pc3}, \texttt{pc4}, \texttt{jm1}) y los de reconocimiento de
        dígitos manuscritos. Las instancias son distintas y la similitud de
        dominio es una propiedad legítima que el recomendador debería poder
        aprovechar; limitarla protege la diversidad del corpus, no su validez.

        Aplicada sobre los 52 candidatos quedan 42 familias distintas, margen
        sobrado para los 15--25 del corpus.

  \item \textbf{Regla de desempate.} Los candidatos se ordenan por IQR
        descendente y se seleccionan los primeros hasta completar el corpus,
        sujeto a dos restricciones de diversidad que se aplican \emph{después}
        del orden:
        \begin{itemize}
          \item cobertura de los tres regímenes de tamaño (centenas, millares y
                decenas de millares de instancias);
          \item al menos cinco datasets con variables categóricas y al menos
                cinco con valores faltantes, para que la codificación y la
                imputación del espacio de búsqueda queden ejercitadas.
        \end{itemize}
        Cuando una restricción obliga a omitir un candidato mejor ubicado en el
        orden, se registra cuál y por qué.

  \item \textbf{Verificación de licencias.} La elegibilidad por licencia se
        comprobó a nivel de suite al fijar los criterios. La licencia individual
        de cada dataset seleccionado se registra junto a su ficha.

  \item \textbf{Umbral del subconjunto sensible.} Un dataset se considera
        sensible al preprocesamiento si su IQR es mayor o igual a $0{,}02$ de
        macro-F1 (dos puntos porcentuales). Este umbral queda fijado aquí,
        \emph{antes} de observar la distribución de los valores, y es el mismo
        que define la estratificación de la Sección~\ref{sec:analisis}.

  \item \textbf{Persistencia.} Los doce valores de cada candidato se conservan en
        el repositorio, incluidos los de los candidatos descartados, para que la
        selección sea verificable y no una afirmación.
\end{enumerate}

\subsection{Dónde queda registrado el corpus resultante}

La identidad de los datasets seleccionados es un \textbf{resultado} de aplicar
este procedimiento, no una regla, de modo que no se registra en este documento
(véase la Sección~\ref{sec:fuera}). Se escribe en
\texttt{docs/experimentos/CORPUS-seleccionado.md}, generado por el mismo script
que ejecuta el cribado, con el IQR de cada candidato, cuántos pipelines se
descartaron por inevaluables, la licencia y el motivo de cada exclusión.

Los datos de partida ya están verificados. Los 52 candidatos se descargaron y sus
metadatos se comprobaron contra los datos reales, sin discrepancias sin explicar;
el detalle está en \texttt{CORPUS-descargado.md}, en la misma carpeta.

\subsection{Resultados posibles declarados de antemano}

Si todos los datasets seleccionados superan el umbral, el subconjunto sensible
coincide con el corpus completo y la estratificación de la
Sección~\ref{sec:analisis} queda vacía. En ese caso se reporta como tal. El
umbral no se redefine después de observar la distribución para forzar una
partición del corpus.

% =========== 6. RECOMENDADOR ================================
\section{El recomendador de arranque en caliente}
\label{sec:recomendador}

El anteproyecto fija dos cosas: que el arranque en caliente inyecta cinco
configuraciones y que el recomendador es de vecinos más cercanos, dejando «un
meta-modelo más sofisticado» como extensión deseable. No especifica cómo se mide
la distancia entre datasets ni cómo se eligen las cinco configuraciones. Esta
sección cierra esos silencios.

\textbf{Todas las reglas siguientes se fijaron a priori}, sin medir sobre el
corpus. La razón es que los 25 datasets del corpus \emph{son} los pliegues del
\textit{leave-one-dataset-out}: elegir un parámetro porque funciona bien sobre
ellos sería seleccionar sobre el conjunto de prueba.

\subsection{Distancia y normalización}

La distancia entre dos datasets es la \textbf{$L_1$ sobre el vector de
meta-características escalado}. La elección de $L_1$ tiene precedente publicado
para exactamente esta tarea: auto-sklearn ordena los datasets por su distancia
$L_1$ en el espacio de meta-características y usa los más cercanos para arrancar
en caliente su optimización bayesiana (Feurer et al., 2015, sección 3.1).

El escalado \textbf{no} proviene de esa referencia, que no describe ninguno, y se
declara aquí como decisión propia: se usan la \textbf{mediana y el rango
intercuartílico}, y no la media y la desviación típica, porque con dos docenas de
datasets un solo valor extremo desplaza ambos estadísticos. Escalar es
imprescindible porque el corpus mezcla meta-características acotadas con otras
que van de 500 a 45\,000; sin escalar, esas últimas dominarían la distancia por su
magnitud.

Como análisis de sensibilidad declarado de antemano, con el mismo patrón que
\texttt{n\_startup\_trials}, se admite repetir el análisis con escalado por media
y desviación y con distancia euclídea. Esa repetición es secundaria y no
sustituye al análisis primario.

\subsection{Cuántos vecinos}

\textbf{El número de vecinos no es un parámetro.} Lo que el anteproyecto fija son
las cinco configuraciones inyectadas, de modo que el recomendador recorre los
vecinos por cercanía y \textbf{salta las configuraciones repetidas} hasta reunir
cinco \emph{distintas}.

\textit{Razón:} inyectar dos veces la misma configuración gastaría dos de los
treinta ensayos del presupuesto para obtener el mismo valor. Con 17
configuraciones ganadoras distintas entre los 25 datasets del corpus, los
repetidos aparecen con frecuencia y no son un caso de borde.

\subsection{Consulta fuera del rango del pliegue}

Antes de escalar, cada valor del vector consultado se \textbf{recorta al mínimo y
al máximo observados en el pliegue}. El recorte afecta únicamente a la consulta,
nunca al ajuste, de modo que no abre ninguna vía de fuga.

\textit{Razón:} el escalado robusto controla los valores extremos \emph{dentro}
del conjunto de ajuste, pero en el \textit{leave-one-dataset-out} el dataset
excluido puede caer muy fuera. Sin recorte, y medido sobre el corpus real, hay dos
pliegues en los que una sola meta-característica aporta más de la mitad de la
distancia —hasta el 71\%—, de modo que el vecino más parecido no lo elegiría un
espacio de 87 meta-características sino una estadística de covarianza. El recorte
no necesita ninguna constante arbitraria: el recomendador no puede medir
\emph{cuánto} fuera del corpus cae un dataset, solo que cae fuera.

\subsection{Meta-características supervivientes, por pliegue}
\label{sec:supervivientes}

La Sección~\ref{sec:metacar} establece que se descarta toda meta-característica no
finita en algún dataset del corpus, y añadía que ese descarte se calcula «una sola
vez, sobre el corpus definitivo». \textbf{Esta versión lo corrige}: el descarte se
recalcula \textbf{en cada pliegue}, usando solo sus datasets de entrenamiento, y
se interseca además con las meta-características finitas del dataset consultado.

\textit{Razón:} calcularlo sobre el corpus completo usaría el patrón de finitud
del dataset excluido para decidir qué columnas existen, que es una forma —débil,
porque solo usa el hecho binario de ser finito y nunca un valor— de fuga en el
nivel meta. La regla por pliegue la elimina.

\textit{Costo:} ninguno. Medido sobre el corpus, solo un pliegue conserva
meta-características que el corpus completo descarta, son dos, y ambas resultan no
finitas en el propio dataset excluido, de modo que la intersección las descarta
igual. Las 88 supervivientes son las mismas en los 25 pliegues bajo cualquiera de
las dos reglas. Se adopta la estricta porque no cuesta nada y elimina la fuga.

\subsection{Las cuatro estrategias y su papel}
\label{sec:estrategias}

Las cuatro estrategias devuelven configuraciones ganadoras de datasets reales de
la base y saltan las repetidas. Lo único que las distingue es \textbf{de qué
datasets} toman las configuraciones, y esa es la propiedad que las hace
comparables: una estrategia que extrajera configuraciones arbitrarias del espacio
mezclaría el efecto de la similitud con el de la calidad de lo inyectado y no
mediría ninguno de los dos.

\begin{center}
\begin{tabularx}{\textwidth}{L{2.2cm} L{4.6cm} Y}
\toprule
\rowcolor{tablehead} \textbf{Estrategia} & \textbf{De dónde toma} & \textbf{Papel en el análisis} \\
\midrule
\texttt{nearest} & Los vecinos más cercanos &
\textbf{Primaria.} Es la que produce el contraste confirmatorio de la
Sección~\ref{sec:analisis}. \\
\midrule
\texttt{random} & Ganadoras de datasets al azar &
\textbf{Ablación, secundaria.} Aísla si lo que aporta es la similitud. \\
\midrule
\texttt{frequent} & Las que más ganan en la base, ignorando las
meta-características &
\textbf{Control, secundario.} Aísla la contribución del meta-aprendizaje. \\
\midrule
\texttt{diverse} & El subconjunto más separado entre las cercanas &
\textbf{Exploratoria.} No puede sostener conclusiones. \\
\bottomrule
\end{tabularx}
\end{center}

\textbf{La ablación aleatoria toma ganadoras de datasets al azar, y no
configuraciones al azar del espacio.} Son dos preguntas distintas: la primera
aísla si la \textbf{similitud} aporta algo, que es lo que el TFG afirma; la
segunda solo preguntaría si arrancar en caliente aporta algo, y al poder caer en
la región inevaluable mezclaría el efecto de la similitud con el de la validez de
las configuraciones. Emplea la semilla $202$.

\textbf{El control \texttt{frequent} se declara porque sin él un resultado
positivo tampoco sería atribuible.} Si el arranque en caliente redujera los
ensayos hasta el objetivo, sin este control no se podría distinguir «los datasets
parecidos comparten buen preprocesamiento» de «hay configuraciones que van bien
casi siempre y conviene probarlas primero». La segunda explicación no necesita
meta-aprendizaje alguno.

% =========== 7. PROTOCOLO LODO ==============================
\section{Protocolo de comparación}
\label{sec:lodo}

El protocolo es el de la Sección~4.3.3 del anteproyecto: validación cruzada por
exclusión de dataset, dos condiciones sobre el dataset excluido con un
presupuesto de 30 ensayos cada una, arranque en frío contra arranque en caliente
con cinco configuraciones inyectadas mediante \texttt{enqueue\_trial}. Lo que
sigue precisa sus puntos no especificados.

\subsection{Semillas}

Cada condición se ejecuta con \textbf{cinco semillas}, y son las \textbf{mismas
cinco} en ambas condiciones: $\{101,\ 202,\ 303,\ 404,\ 505\}$. La construcción de
la base de meta-conocimiento, que el anteproyecto define con un presupuesto único
de 50 ensayos por dataset, emplea la semilla $42$. El cribado de sensibilidad
emplea la semilla $4242$. La estrategia de recomendación aleatoria de la
Sección~\ref{sec:recomendador} emplea la semilla $202$.

Repetir con varias semillas responde a que una sola corrida por condición produce
una observación pareada ruidosa: el muestreador TPE con 30 ensayos tiene varianza
apreciable entre semillas, y una conclusión que dependa de la semilla que tocó no
es defendible. Esto no altera el presupuesto de 30 ensayos, ni el emparejamiento
por dataset, ni el número de pares de la prueba estadística.

\textbf{Decisión cerrada en esta versión: se pasa de tres semillas a cinco.} La
v0.1 fijó tres y condicionó el ascenso a cinco al costo de cómputo medido, que la
v0.2 aplazó por no existir todavía. La medición ya existe y descarta el cómputo
como obstáculo: la base de meta-conocimiento completa, con 50 ensayos sobre cada
uno de los 25 datasets, se construyó en 1,1 minutos, y el cribado de los 52
candidatos tardó 42 segundos. A ese orden de magnitud, el
\textit{leave-one-dataset-out} completo son minutos con tres semillas y siguen
siendo minutos con cinco.

Las tres semillas originales se conservan y se añaden dos: la decisión amplía y
no sustituye. El motivo de ampliar es que \textbf{cinco observaciones por dataset hacen
más estable la mediana} de la que depende cada par de la prueba de Wilcoxon, y
que con la regla de censura de la Sección~\ref{sec:censura} hace falta que la
mayoría de las semillas de un dataset queden censuradas para desplazar su
observación: con tres bastaban dos, con cinco hacen falta tres.

En la implementación, la semilla gobierna tanto la partición de entrenamiento y
validación como el muestreador de Optuna. Como ambas condiciones comparten la
semilla, dentro de cada semilla la comparación entre frío y caliente se realiza
sobre una partición idéntica.

\subsection{Definición del objetivo}

El objetivo se define por pareja \textbf{(dataset, semilla)}: para el dataset $d$
y la semilla $s$, el objetivo es el 95\% del mejor valor de macro-F1 que
alcanza el arranque en frío con la semilla $s$ sobre el dataset $d$ en su
presupuesto completo de 30 ensayos. Los ensayos hasta el objetivo se calculan por
semilla, y la observación pareada de cada dataset es la \textbf{mediana de las
cinco semillas} en cada condición.

\textit{Razón:} la alternativa —definir un objetivo único por dataset a partir de
la agregación de las cinco corridas en frío— permitiría que el arranque en frío no
alcanzara su propio objetivo en alguna semilla, reintroduciendo la censura en la
condición de referencia. La definición por pareja preserva el emparejamiento de
semillas que exige el anteproyecto y mantiene la propiedad de la
Sección~\ref{sec:censura}.

\subsection{Los ensayos inyectados cuentan}

Las cinco configuraciones inyectadas como arranque en caliente cuentan dentro de
los ensayos hasta el objetivo. El arranque en caliente dispone por tanto de
veinticinco ensayos de exploración propia dentro del presupuesto de treinta.

\textit{Razón:} no contarlas otorgaría al arranque en caliente una ventaja
puramente contable, ajena a la hipótesis, que se investiga sobre el costo total
de la búsqueda.

\subsection{Cómo se cuenta la variable}
\label{sec:conteo}

Estas tres precisiones no cambian el protocolo: cierran silencios que solo se
hacen visibles al implementar el arnés de medición, y quedan aquí registradas
antes de ejecutarlo.

\begin{enumerate}
  \item \textbf{El conteo empieza en uno.} «Ensayos hasta el objetivo» es la
        posición del primer ensayo cuya métrica alcanza el umbral, con el primer
        ensayo de la búsqueda numerado como $1$. Es lo coherente con que la regla
        de censura de la Sección~\ref{sec:censura} valga $31$, es decir
        presupuesto más uno.

  \item \textbf{Los ensayos inevaluables consumen su turno.} Un candidato que no
        se puede evaluar sobre el dataset se registra con métrica \texttt{NaN},
        nunca alcanza el objetivo y \textbf{no se descuenta} del conteo. Vale
        tanto para el arranque en frío, que llega a esa región por muestreo, como
        para el caliente, cuyas configuraciones provienen de otros datasets y
        pueden no ser ejecutables en el receptor.

  \item \textbf{Se registra la curva completa por ensayo.} De cada corrida se
        guardan las treinta métricas en orden, y no solo el ensayo en que se
        cruza el umbral. Es lo que permite calcular la medida descriptiva de
        desempeño a quince ensayos que define la Sección~4.3.3 del anteproyecto,
        y recalcular la variable bajo las reglas de censura alternativas de la
        Sección~\ref{sec:censura} sin volver a ejecutar nada.
\end{enumerate}

\subsection{Casos censurados}
\label{sec:censura}

Cuando una condición no alcanza el objetivo dentro del presupuesto, su valor de
ensayos hasta el objetivo se fija en $31$ (presupuesto más uno). La regla se
aplica de forma simétrica a ambas condiciones y el número de casos censurados se
reporta por condición y a nivel de semilla, no de dataset.

\textbf{La censura es asimétrica por construcción, y conviene enunciarlo con
precisión: el arranque en frío no puede censurarse.} El objetivo es el 95\% del
mejor valor que el propio arranque en frío alcanza en su presupuesto completo; si
ese mejor valor ocurre en el ensayo $t \leq 30$, entonces el frío cruzó el 95\%
en algún ensayo menor o igual que $t$. No se trata de un suceso improbable sino
imposible. Por lo tanto la regla de censura solo puede afectar al arranque en
caliente, y su elección requiere una garantía explícita, que es la siguiente.

\textbf{Chequeo de robustez.} El análisis primario se recalcula íntegro con
castigos más severos para los casos censurados: dos veces y diez veces el
presupuesto ($60$ y $300$). Si la conclusión no cambia con ninguna de las tres
reglas, se declara explícitamente en el informe que la conclusión no depende de
la regla de censura. Si cambia con alguna, se reporta con qué regla y en qué
sentido, sin seleccionar la más favorable.

Con la mediana de cinco semillas, la regla de censura solo desplaza la observación
de un dataset cuando al menos tres de sus cinco semillas quedan censuradas; el
reporte a nivel de semilla hace visible esa situación.

\subsection{Ausencia de poda efectiva}

El objetivo evalúa una única partición y no reporta valores intermedios, de modo
que el podador por mediana configurado en el motor de búsqueda no interviene.
Los treinta ensayos son treinta evaluaciones completas y ninguno queda sin valor
de métrica, por lo que no se requiere una regla para ensayos podados.

\subsection{Ausencia de fuga en el nivel meta}

La normalización del espacio de meta-características y la métrica de distancia
del recomendador se ajustan \textbf{únicamente con los datasets de entrenamiento
de cada pliegue}, nunca con el dataset excluido. Ajustarlas sobre el corpus
completo contaminaría cada pliegue y inflaría la ventaja del arranque en caliente
de un modo difícil de detectar a posteriori. Es el análogo, en el nivel meta, del
ajuste de las transformaciones con estado usando solo el subconjunto de
entrenamiento.

% =========== 6. PLAN DE ANALISIS ============================
\section{Plan de análisis}
\label{sec:analisis}

\subsection{Análisis primario}

Sobre el \textbf{corpus completo}, se contrastan las medianas de ensayos hasta el
objetivo de ambas condiciones. La hipótesis se considera confirmada si se cumplen
\emph{ambas} condiciones de la Sección~4.3.4 del anteproyecto:

\begin{enumerate}
  \item la mediana de ensayos hasta el objetivo del arranque en caliente es al
        menos un 30\% menor que la del arranque en frío; \emph{y}
  \item la prueba de rangos con signo de Wilcoxon, pareada por dataset, resulta
        significativa con $\alpha = 0{,}05$.
\end{enumerate}

La prueba se realiza \textbf{de una cola}, en la dirección de la hipótesis
(reducción), y se reporta además el valor $p$ de dos colas. La lateralidad se fija
aquí, antes de ver los datos: la hipótesis del anteproyecto es direccional
(«reduce la mediana en al menos un 30\%»), lo que justifica la prueba de una
cola, y con $N$ entre 15 y 25 pares la potencia disponible es limitada. Reportar
ambos valores evita que la elección de lateralidad se convierta en un grado de
libertad oculto.

La variable «desempeño a presupuesto reducido» (mejor valor de métrica tras
quince ensayos), definida en la Sección~4.3.3 del anteproyecto, se reporta como
medida descriptiva. No forma parte del criterio de confirmación, que la
Sección~4.3.4 hace depender exclusivamente de los ensayos hasta el objetivo.

\subsection{Análisis secundarios, declarados como tales}

\begin{enumerate}
  \item \textbf{Estratificado por sensibilidad.} La misma prueba sobre el
        subconjunto de datasets con IQR $\geq 0{,}02$, etiquetada como secundaria
        y acompañada de su $N$. Si el subconjunto queda por debajo de diez
        datasets, se reportan el tamaño del efecto y las medianas sin
        sobreinterpretar el valor $p$.
  \item \textbf{Robustez frente a la regla de censura}, según la
        Sección~\ref{sec:censura}.
  \item \textbf{Tabla por dataset}, con sensibilidad, margen sobre el pipeline
        vacío, ensayos hasta el objetivo por semilla y condición, casos
        censurados, desempeño a quince ensayos y las dos covariables de la
        Sección~\ref{sec:covariables}. Los datasets planos aparecen en su propia
        fila con su valor de sensibilidad visible, no promediados dentro de la
        mediana global.
  \item \textbf{Ablación aleatoria.} El mismo contraste sustituyendo la
        estrategia primaria por \texttt{random} (Sección~\ref{sec:estrategias}).
        Responde a si lo que aporta el arranque en caliente es la
        \textbf{similitud} o el mero hecho de inyectar ganadoras plausibles.
  \item \textbf{Control por frecuencia.} El mismo contraste sustituyendo la
        estrategia primaria por la que recomienda, sin mirar las
        meta-características, las configuraciones que más ganan en la base.
        Responde a si el meta-aprendizaje aporta algo sobre esa regla.
\end{enumerate}

Ningún análisis secundario sustituye al primario en la conclusión del proyecto.
Los dos últimos existen porque \textbf{condicionan cómo se puede interpretar un
resultado positivo}, no solo uno negativo: si el arranque en caliente redujera los
ensayos hasta el objetivo pero \texttt{frequent} lograra lo mismo, la reducción no
sería atribuible al meta-aprendizaje.

\subsection{Covariables declaradas}
\label{sec:covariables}

Se registran de antemano dos covariables descriptivas, ambas \textbf{ya medidas
antes de escribir esta versión} y ninguna derivada del experimento. No entran en
el criterio de confirmación de la Sección~4.3.4 del anteproyecto: su función es
que el resultado, sea cual sea, pueda interpretarse.

\begin{enumerate}
  \item \textbf{Fracción de ensayos inevaluables}, por dataset y condición. Hace
        visible dataset por dataset la asimetría que la
        Sección~\ref{sec:espacio} declara y no corrige.

  \item \textbf{Anchura de la meseta}: cuántas de las 144 configuraciones del
        espacio alcanzan el 95\% del óptimo de ese dataset. Se mide por rejilla
        exhaustiva y queda registrada en el documento
        % \allowbreak permite partir el nombre del archivo sin añadir un guion:
        % dentro de la sangría de la lista no cabe entero en una línea.
        \texttt{DIVERSIDAD-\allowbreak ganadores.md} de esta misma carpeta.
\end{enumerate}

\textbf{Por qué la segunda, y por qué antes de ejecutar.} La medición previa da
una mediana de 36 de 144 configuraciones dentro de la meseta, lo que equivale a
unos cuatro ensayos aleatorios esperados hasta el objetivo en el dataset mediano,
y menos de dos en cuatro datasets. Donde la meseta es ancha, el arranque en frío
cruza el objetivo dentro de su propia fase aleatoria y no queda margen para
mejorarlo: el contraste no puede discriminar por construcción, con independencia
de la calidad del recomendador.

Declararla ahora tiene una consecuencia concreta sobre lo que el proyecto podrá
afirmar. Sin ella, un resultado nulo solo permitiría decir que no se observó
reducción. Con ella, permite decir \emph{en qué datasets había margen para
observarla y en cuáles no}, que es una afirmación sustantiva sobre cuándo puede
ayudar el arranque en caliente meta-aprendido y no una excusa construida después
de ver los resultados. Por eso se declara antes y no después.

\subsection{Análisis exploratorios, no confirmatorios}

Se declaran exploratorios, y por tanto no pueden usarse para sostener la
hipótesis:

\begin{itemize}
  \item el análisis de diversidad de los pipelines ganadores de la base de
        meta-conocimiento (semana 7);
  \item la ablación con un recomendador aleatorio como control (semana 11);
  \item cualquier análisis de causas posterior a la observación de los
        resultados.
\end{itemize}

Su valor es interpretativo y así se presentan en el informe.

\subsection{Resultado negativo}

Conforme a la mitigación del riesgo RI1 del anteproyecto, si el criterio de la
Sección~4.3.4 no se alcanza, el resultado se reporta como tal, sin reformular la
hipótesis ni el umbral. El análisis de causas y la generalización del framework
constituyen resultados defendibles por sí mismos.

% =========== 7. FUERA DE ALCANCE ============================
\section{Qué queda fuera de este documento}
\label{sec:fuera}

\begin{itemize}
  \item La identidad de los datasets del corpus y sus valores de IQR: son
        resultado de aplicar la Sección~\ref{sec:corpus}, no reglas, y se
        registran en \texttt{docs/experimentos/CORPUS-seleccionado.md}. La v0.1
        anunciaba que irían en la v0.2, lo que contradice el último punto de esta
        lista y la regla de que este documento se cierra antes de ejecutar.
  \item La lista concreta de meta-características descartadas por no finitas: es
        el resultado de aplicar la regla de la Sección~\ref{sec:metacar} al corpus
        definitivo, y se registra junto a la base de meta-conocimiento.
  \item La incorporación o el descarte del agrupamiento como extensión deseable:
        se decidió en el punto de control de alcance de la semana 7 y quedó
        \textbf{fuera}, sin afectar al protocolo aquí registrado
        (\texttt{docs/decisiones/ADR-003-alcance-agrupamiento.md}).
  \item Todo resultado experimental. Este documento registra reglas, no
        hallazgos.
\end{itemize}

\subsection{Reglas que gobiernan entradas ya construidas}

Una precisión necesaria para que nadie tenga que adivinarla. Dos de las reglas
registradas en esta versión \textbf{ya se habían aplicado} cuando se escribió:

\begin{itemize}
  \item El \textbf{desempate canónico} entre configuraciones ganadoras
        empatadas, que gobernó la construcción de la base de meta-conocimiento en
        la semana 7.
  \item Las reglas del \textbf{recomendador} de la
        Sección~\ref{sec:recomendador}, aplicadas en la evaluación exploratoria
        sobre datasets fuera del corpus de la semana 8.
\end{itemize}

Ninguna de las dos vulnera la regla de oro, porque esa regla protege el
\textbf{experimento} y ninguna de esas dos ejecuciones es un contraste de la
hipótesis: la base es una entrada del experimento y la evaluación fue exploratoria
y se declaró como tal, sin usarse para elegir ningún parámetro. Pero conviene que
el lector pueda distinguir sin esfuerzo qué se fijó antes de qué, así que se dice
aquí en lugar de dejarlo implícito.

Todo lo demás que esta versión registra gobierna el \textit{leave-one-dataset-out},
que no se había ejecutado ni siquiera en parte cuando el documento se cerró.

% =========== 8. HISTORIAL ===================================
\section{Historial de versiones}

\begin{center}
\begin{tabularx}{\textwidth}{L{1.6cm} L{2.4cm} Y}
\toprule
\rowcolor{tablehead} \textbf{Versión} & \textbf{Fecha} & \textbf{Contenido y motivo} \\
\midrule
v0.1 & \preregfechauno &
Versión inicial. Fija el modelo base, la parametrización relativa y el versionado
del espacio de búsqueda, el procedimiento de cribado de sensibilidad, las
semillas, la definición del objetivo, el conteo de los ensayos inyectados, la
regla de censura con su chequeo de robustez y el plan de análisis. Se cierra en
la semana 1 para que las decisiones queden fechadas el día en que se tomaron. \\
\midrule
v0.2 & \preregfechados &
Congela la composición del espacio de búsqueda \texttt{tabular-v1} y declara su
región inevaluable; fija las reglas de extracción de meta-características
(biblioteca, grupos, submuestreo de 500 con su semilla y estratificación,
imputación previa, redondeo y descarte de no finitas); añade la regla de familia y
el remuestreo de los pipelines inevaluables en el cribado; y precisa dónde queda
registrado el corpus resultante. Se cierra \textbf{antes} de ejecutar el cribado
de sensibilidad. Corrige además dos anuncios de la v0.1 que eran irrealizables:
el costo de cómputo y la identidad del corpus no pueden registrarse en esta
versión (véase la Sección~\ref{sec:fuera}). Lleva una errata sobre su fechado,
registrada en la Sección~\ref{sec:erratas}. \\
\midrule
v1.0 & \preregfecha &
\textbf{Versión congelada previa a la ejecución del experimento.} Cierra las dos
decisiones que quedaban abiertas: el número de semillas asciende de tres a cinco,
porque el costo medido descarta el cómputo como obstáculo, y la región inevaluable
se deja como está en lugar de introducir pasos obligatorios, declarándola en su
lugar como covariable. Añade la Sección~\ref{sec:recomendador} con las reglas del
recomendador —distancia, escalado, número de vecinos, recorte de la consulta y
supervivientes por pliegue—, que corrige en ese último punto lo que la v0.2 decía;
la Sección~\ref{sec:estrategias} con las cuatro estrategias y su papel en el
análisis; la Sección~\ref{sec:conteo} con las tres precisiones sobre cómo se
cuenta la variable medida; y la Sección~\ref{sec:covariables} con las dos
covariables declaradas. Se cierra \textbf{antes} de generar cualquier resultado
experimental. \\
\bottomrule
\end{tabularx}
\end{center}

% =========== 9. FE DE ERRATAS ===============================
\section{Fe de erratas}
\label{sec:erratas}

Las correcciones posteriores al cierre de una versión se registran aquí, en lugar
de aplicarse en silencio sobre el texto congelado. \textbf{Ninguna de ellas puede
modificar una regla del protocolo}: si una regla tuviera que cambiar, eso exige una
versión nueva y no una errata.

El estado exacto de cada versión en el momento de cerrarse queda preservado por su
etiqueta y se recupera con \texttt{git show \preregtag}, de modo que cualquier
diferencia entre este documento y el texto etiquetado es auditable.

\subsection*{2026-09-04 --- Fechado de la versión v0.2}

\textbf{Qué decía.} La portada afirmaba que la etiqueta \texttt{prereg-v0.2} era
una «etiqueta anotada creada al cerrar esta versión» y que «su fecha de creación,
registrada por git, es el fechado verificable de este documento».

\textbf{Por qué es incorrecto.} La etiqueta se creó el 4 de septiembre de 2026,
siete días después del commit que cierra la versión, y no el mismo día como sí
ocurrió con la v0.1. Su fecha de creación es por tanto posterior a las ejecuciones
que esta versión gobierna, y no sirve como fechado de nada.

\textbf{Corrección.} La portada pasa a señalar como fechado verificable la fecha
del \emph{commit} etiquetado, que sí es anterior a esas ejecuciones. La etiqueta
conserva su función de identificar la versión, pero no la de fecharla.

\textbf{Qué sostiene el fechado.} El orden de los commits, que es independiente de
la etiqueta y no puede alterarse sin que el historial lo delate:

\begin{center}
\begin{tabularx}{0.92\textwidth}{L{2.2cm} L{3.4cm} Y}
\toprule
\rowcolor{tablehead} \textbf{Commit} & \textbf{Fecha} & \textbf{Qué hizo} \\
\midrule
\texttt{8e4cb10} & 2026-08-28 11:10 & Cierra el pre-registro v0.2 \textit{(commit etiquetado)} \\
\texttt{014da98} & 2026-08-28 11:35 & Ejecuta el cribado y selecciona el corpus \\
\texttt{18fc1a3} & 2026-08-28 15:36 & Construye la base de meta-conocimiento \\
\bottomrule
\end{tabularx}
\end{center}

El documento se cerró veinticinco minutos antes del cribado y cinco horas antes de
la construcción de la base, de modo que la regla de oro de la Sección~1 se cumple y
es comprobable por cualquiera con acceso al repositorio.

\textbf{Qué no se hizo.} No se retrodató la etiqueta. Git permite fijar su fecha de
creación mediante variables de entorno, y hacerlo habría producido un fechado
aparentemente correcto y falso, que es exactamente lo que un pre-registro existe
para impedir.

\textbf{Alcance de esta errata.} Corrige una afirmación sobre el \emph{mecanismo}
de fechado. No altera ninguna decisión de protocolo, ningún criterio de análisis ni
ninguna regla registrada en las secciones anteriores, y no afecta a ningún
resultado ya producido bajo esta versión.

\textbf{Consecuencia para la v1.0.} Esa versión precede a la ejecución del
experimento, así que su etiqueta debe crearse \textbf{el mismo día} que el commit
que la cierra, como se hizo con la v0.1. De lo contrario volvería a hacer falta una
errata como esta, en la versión donde menos margen hay para ello.

\subsection*{2026-09-04 --- Composición de la configuración del modelo base}

\textbf{Qué pasaba.} En el PDF de las versiones v0.1 y v0.2, la configuración
congelada del modelo base de la Sección~\ref{sec:modelo} se imprimía como
\texttt{weights=üniform"} en lugar de \texttt{weights="uniform"}.

\textbf{Causa.} Con \texttt{babel} en español la comilla recta es un carácter
activo, y la secuencia \texttt{"u} compone «ü». Afectaba solo a ese parámetro:
\texttt{metric="minkowski"} y \texttt{algorithm="brute"} salían bien, porque
\texttt{"m} y \texttt{"b} no son atajos.

\textbf{Corrección.} Se desactiva el atajo alrededor de ese bloque. El código
fuente del documento \emph{siempre} dijo \texttt{weights="uniform"}: lo que estaba
mal era la composición tipográfica, no el valor registrado.

\textbf{Alcance.} No cambia el modelo base ni ninguna decisión de protocolo. Se
registra porque el literal afectado es el que hay que copiar para reproducir el
experimento, y quien lo tomara del PDF de una versión anterior obtendría un
parámetro que scikit-learn rechaza.

\end{document}
